% Einheit 2 von 4 – Gleichungen mit e-Funktion und Logarithmus (bank/gleichungen-loesen/e2.jsonl, zusammenbau v0.1)
%% TODO Zeitmarke Einheit 2: Marken-Zeile nicht lesbar
%% TODO Zweigzeile Teil 1: was hier gelernt wird (Satz in Schülersprache aus dem Einheitstitel) – Einheit 2
\einheitenkopf[e2]{Einheit 2 von 4 · Gleichungen mit e-Funktion und Logarithmus}
\zweigzeile{Abitur GK}
\setcounter{aufgabe}{13}

%% TODO Titel: Ich-kann-Satz fehlt in der Bank (Platzhalter: Kettenname) – Nr. 14
%% TODO Anweisung: ein Satz über der ersten Teilaufgabe fehlt in der Bank – Nr. 14
\begin{aufgabe}{Gleichungen mit e-Funktion und Logarithmus}
%% TODO \rechenplatz{n}: Schrittzahl des Grundfalls fehlt in der Bank (nur bei fester Schreibform, 2.3 b)
\begin{teile}
\teil Welcher Weg? $4e^{3x} - 20 = 0$ \kreuz{Logarithmieren} \\ \kreuz{e-Faktor kürzen} \\ \kreuz{e-Funktion anwenden}
\end{teile}
\begin{gleichungsraster}[2]
\gl{\text{Löse: $4e^{2x} - 12 = 0$. Runde auf zwei Stellen.}} &
\gl{\text{$f(x) = 3e^{2x} - 1$. Bestimme die Nullstelle exakt. (Abitur 2017 LK)}} \\
\gl{\text{Die Konzentration eines Wirkstoffs im Blut ist $c(t) = 80 \cdot e^{-0{,}2t}$ ($t$ in h, $c$ in mg/l). Anfangskonzentration? Wann sind nur noch 10 mg/l vorhanden? Runde auf zwei Stellen. (Abitur 2022 LK)}} &
\gl{\text{Die Zahl der Biber in einem Revier ist $B(t) = \frac{60}{1 + 5e^{-0{,}5t}}$ ($t$ in Jahren). Anzahl zu Beginn? Wann sind es 50 Biber? Runde auf zwei Stellen. (Abitur 2024 LK)}} \\
\gl{\text{Eine Rutsche hat das Profil $f(x) = 0{,}5 + 2e^{0{,}4x}$ für $-3 \le x \le 0$ (1 LE = 1 m, $x = -3$ ist das untere Ende). Wie weit ist die Stelle in 2 m Höhe waagerecht vom unteren Ende entfernt? Runde auf zwei Stellen. (Abitur 2018 LK)}} &
\gl{\text{$f(x) = (x^2 - 3x) \cdot e^{-0{,}5x}$, $g(x) = x \cdot e^{-0{,}5x}$ (1 LE = 1 dm). Wie weit liegen die Schnittstellen waagerecht auseinander? (Abitur 2021 GK)}} \\
\end{gleichungsraster}
\begin{teile}
\teil $f(x) = 9e^{x - 2}$, $g(x) = x^2 \cdot e^{x - 2}$. Zeige, dass sich die Graphen nur für $x = -3$ und $x = 3$ schneiden. \hfill (Abitur 2026 LK)
\teil $f(x) = (x - 1)^2 \cdot e^{x}$ hat die Ableitung $f'(x) = (x^2 - 1) \cdot e^{x}$. Zeige, dass sich die Graphen von $f$ und $f'$ an der Stelle 1 schneiden. \hfill (Abitur 2025 LK)
\teil $f(x) = e^{x} + 2x$, $g(x) = 2x - 3$. Begründe, dass die Graphen keinen gemeinsamen Punkt haben. \hfill (Abitur 2018 LK)
\end{teile}
\begin{gleichungsraster}[2]
\gl{\text{$f(x) = \mathrm{ln}(x + 4) + 2$. Bestimme die Nullstelle. (Abitur 2026 LK)}} &
\gl{\text{$f(x) = e^{0{,}5x^2}$ hat die Ableitung $f'(x) = x \cdot f(x)$. Die Graphen von $f$ und $f'$ schneiden sich in genau einem Punkt. Steigung des Graphen von $f$ in diesem Punkt? (Abitur 2023 GK)}} \\
\end{gleichungsraster}
\end{aufgabe}

%% TODO Titel: Ich-kann-Satz fehlt in der Bank (Platzhalter: Kettenname) – Nr. 15
%% TODO Anweisung: ein Satz über der ersten Teilaufgabe fehlt in der Bank – Nr. 15
\begin{aufgabe}{Schnittpunkte zweier Graphen durch Ausklammern eines gemeinsamen Faktors berechnen}
\begin{gleichungsraster}[2]
\gl{\text{$f(x) = (x - 2) \cdot e^{0{,}25x}$, $g(x) = x - 2$. Koordinaten der Schnittpunkte? (Abitur 2018 GK)}} & \\
\end{gleichungsraster}
\end{aufgabe}

%% TODO Titel: Ich-kann-Satz fehlt in der Bank (Platzhalter: Kettenname) – Nr. 16
%% TODO Anweisung: ein Satz über der ersten Teilaufgabe fehlt in der Bank – Nr. 16
\begin{aufgabe}{Gleichungen mit e-Funktion und Logarithmus – Fehler finden, Begründen, Anwendung}
\begin{teile}
\teil Finde den Fehler und rechne richtig. \rechnung{3e^{2x} &= 12 &&\mid \mathrm{ln} \\ 2x &= \mathrm{ln}(12) \\ x &\approx 1{,}24}
\teil Begründe, warum man bei $(x - 4) \cdot e^{3x} = 2 \cdot e^{3x}$ durch $e^{3x}$ teilen darf.
\end{teile}
\begin{gleichungsraster}[2]
\gl{\text{Ein Kaffee kühlt nach $T(t) = 20 + 65 \cdot e^{-0{,}05t}$ ab ($t$ in min, $T$ in °C). Ab 60 °C kann man ihn trinken. Nach wie vielen Minuten ist das? Runde auf zwei Stellen.}} & \\
\end{gleichungsraster}
\end{aufgabe}
